\section{Related work}
\label{sec:related}

\textbf{Recursive self-improvement is the reason this question is asked at all.} The premise that AI
systems will soon carry out AI research is explicit in forecasts and in laboratory practice, and it has
produced a line of work that treats machine-learning engineering as its executable testbed --- what
Frontis-MA1 calls AI4AI, post-training a meta-evolution model whose search is composed from atomic
program-evolution operators \citep{yang2026frontis}; RSIBench-Data, which fixes the post-training stack
so that only the agent's \emph{research} decisions vary, and finds that agents improve on their first valid
attempt in $58.3\%$ of settings by refining training-data strategies from feedback, yet that among
searches continuing past their peak, $78.3\%$ finish below it \citep{meng2026rsibench}; and the Meta-Agent Challenge, which asks a model to
build an agent rather than a recipe and reports that meta-agents rarely match human-engineered
baselines \citep{lu2026metaagent}. Around it sits a broader effort to make AI research itself an evaluable
activity: PaperBench asks agents to replicate twenty ICML papers from scratch against author-built
rubrics, where the best tested agent reaches a $21.0\%$ replication score \citep{starace2025paperbench};
EXP-Bench extracts $461$ experimental tasks from $51$ papers and finds agents able to design procedures
far more often than to execute them correctly \citep{kon2025expbench}; MLR-Bench scores $201$
workshop-derived research tasks and reports that coding agents fabricate or fail to validate
experimental results in around $80\%$ of cases \citep{chen2025mlrbench}; and MARS attacks the same
bottleneck from the systems side, arguing that machine-learning engineering is distinctive because
evaluation is expensive and performance attribution opaque \citep{chen2026mars}. Agent$^2$RL-Bench is
closest to our agentic-RL task, testing whether agents can close an interactive RL loop rather than
emit a training script \citep{chen2026agent2rlbench}. We share this line's object of study and its
motivation. What we add is an apparatus in which the improvement is a source-only patch replayed from
a pinned commit, so that a reported gain is an artifact a third party can rebuild rather than a number
produced inside the agent's own workspace.

\textbf{The critique of verifiable-metric benchmarks is the one we take most seriously.}
\citet{kirgis2026openended} argue that the evidence base for autonomous AI R\&D is thin and that what
unifies it --- from Kaggle-derived engineering suites \citep{chan2024mlebench,qiang2025mledojo} to
repository-scale ML tasks \citep{tang2023mlbench} and human-calibrated task batteries
\citep{rein2025hcast} --- is a focus on verifiable tasks with a fixed narrow metric, which cannot reach the parts of
research that matter --- choosing candidate hypotheses, deciding what evidence would settle a question,
and recognising that an approach has failed and that the right move is to start over --- and answer it
by shadow-evaluating open-ended questions against their original authors. Our tasks are verifiable-metric
tasks and inherit that limitation at the level of the \emph{score}. Our response is different from
theirs and complementary to it: because every configuration also leaves an audited trajectory, the same
runs can be read for exactly those judgements without changing the task format
(\S\ref{sec:behavior}) --- which hypothesis a system commits to and when, and whether it ever abandons
the method it was handed. On the tasks where that choice is legible, it separates the artifacts more
cleanly than the reasoning-effort budget does. We read this as evidence that the judgement gap
\citet{kirgis2026openended} identify is measurable inside verifiable tasks, provided the trajectory is
retained and audited, rather than only outside them.

Three further lines of work bear directly on what this paper measures.

\textbf{Benchmarks where agents improve AI systems} give us our closest points of comparison. Two earlier suites already put agents in front of an existing
baseline and asked them to raise it: MLAgentBench, whose thirteen tasks run from CIFAR-10 to BabyLM and
whose best agent averages a $37.5\%$ success rate \citep{huang2024mlagentbench}, and MLGym, whose
thirteen open-ended research tasks yield the finding closest to our own --- frontier models ``improve on
the given baselines, usually by finding better hyperparameters, but do not generate novel'' ideas or
algorithms \citep{nathani2025mlgym}. Our capability result reproduces that shape under a lifecycle that
makes each improvement independently replayable. MLS-Bench is the closest in aim, with 140 tasks in
which an agent improves one component of an ML system and then shows the improvement generalizes
\citep{lyu2026mlsbench}; its findings are ones our data support from a
narrower base --- agents remain far from reliably surpassing human-designed methods, tuning is easier
for them than method invention, and more search does not remove the bottleneck. PostTrainBench is the
closest in protocol, giving ten hours on one H100 to post-train a base model against a released
instruction-tuned checkpoint, where the best agent reaches $23.2\%$ of the official release's $51.1\%$
\citep{rank2026posttrainbench}. That is a score ratio rather than a win rate, so it is not directly
comparable with ours; the point of contact is the reference, not the number. A repository's stock
configuration is a weaker bar than an official release, which is exactly why we also score against a
budget-matched replay of that configuration (\S\ref{sec:validity}). Outside AI systems, Frontier-Eng poses
the same shape of problem --- an existing feasible design, an executable verifier, a fixed interaction
budget --- across 47 industrial engineering tasks, and reports that improvements decay as a dual power law in both frequency and
magnitude, and that search depth matters more than width under a fixed budget
\citep{einsia2026frontiereng}. The decay is the part our data speaks to: buying more search per
configuration lengthens the trajectory without moving the win rate (\S\ref{sec:effort}). The
depth-versus-width finding is about allocating a fixed budget, which our sweep does not vary --- we
change the total, not its distribution --- so we take it as a neighbouring question rather than a
confirmation.

\textbf{Reference-normalized scoring} is what we considered and rejected. MLRC-Bench scores the fraction
of the baseline-to-human gap an agent closes, with its best agent closing $9.3\%$
\citep{zhang2025mlrcbench}. That construction needs a reference distance that is stable and comparable
across tasks, and on this suite it is neither: the distance is missing on two tasks, negatively signed on
one, and not a scalar on a fourth (Table~\ref{tab:tasks}), and where it exists it is set by wherever each
repository's recipe happened to land. We report per-task rates instead (\S\ref{sec:setup}). RE-Bench
supplies the human calibration the field otherwise lacks --- 71 eight-hour expert attempts, with agents
ahead at two hours and humans ahead at eight \citep{wijk2024rebench} --- and that crossover is why we
treat our four-hour/twelve-hour split as a variable rather than a constant. MLE-bench established the
threshold-scored format over 75 Kaggle competitions with human baselines taken from public leaderboards
\citep{chan2024mlebench}, from which we inherit the discipline of scoring against a measured pool
rather than an asserted target. PostTrainBench is also where the failure modes we designed against were
first reported at this scale: agents training on the test set, or downloading an existing
instruction-tuned checkpoint instead of training their own \citep{rank2026posttrainbench}. Our
source-only patch and fresh replay exist to make both mechanically impossible rather than merely
discouraged, and \S\ref{sec:patches} reports the audit that checks it.

\textbf{Validity audits of agent benchmarks} are the line our own metric audit joins. BenchJack
derives a taxonomy of eight recurring flaw patterns from past reward-hacking incidents, then
red-teams ten agent benchmarks and surfaces $219$ distinct flaws across those classes --- synthesizing
exploits that score near-perfectly without solving a single task \citep{wang2026benchjack}; an
audit of $2{,}385$ traces across 15 benchmarks finds exposure or reward hacking in $67.0\%$ of Frontier
Science traces \citep{shao2026doagentbench}; and EvilGenie observes explicit reward hacking by both Codex
and Claude Code \citep{gabor2025evilgenie}, the two harnesses this paper runs, which is why we treated
harness-level exploitation as a live hypothesis rather than a remote one. The measurement strategies
in this line are complementary to ours and worth separating. Some embed exploitable opportunities in
the environment so that taking them is verifiable by construction
\citep{roth2026hackverifiable,thaman2026rhb}; some make the compromise vectors explicit and detect
them by comparing the agent's reported metric against a trusted recomputation
\citep{atinafu2026rha}; one shows that with a test oracle in the loop a coding agent will build to
that oracle and leave the requested artifact hollow \citep{ma2026buildingtotest}; and a survey now
organises the mechanisms across the alignment pipeline \citep{wang2026rhsurvey}. Two further
observations shaped this paper directly: that a saturated benchmark still measures construct validity,
efficiency and reliability if one looks past accuracy \citep{nadgir2026lifeafter}, which is close to
our argument for reading trajectories; and that when the deliverable is an artifact to be deployed,
monitoring must be able to probe the artifact itself \citep{libon2026researcharena}, which is what our
frozen evaluator does by construction. What this paper adds is that our
two qualified tasks (\S\ref{sec:validity}) need different remedies, and only one is about reward hacking:
one is a detection-and-enforcement gap, where an advisory mechanism fired and was ignored by construction;
the other is a specification gap no anti-hacking defense would catch, because nothing was hacked --- an
outcome-defined task admitted a method answering a different question from the one it was written to ask,
and every provenance boundary held while it happened.
